\providecommand{\kBKM}{\kappa_{\mathrm{BKM}}}
\providecommand{\kch}{\kappa_{\mathrm{ch}}}
\providecommand{\kcat}{\kappa_{\mathrm{cat}}}
\providecommand{\kfib}{\kappa_{\mathrm{fiber}}}
\providecommand{\CYcat}{\mathrm{CY}\text{-cat}}
\providecommand{\ChirAlg}{\mathrm{ChirAlg}}
\providecommand{\HolFA}{\mathrm{HolFA}}
\providecommand{\Sp}{\mathrm{Sp}}
\providecommand{\SpCh}{\mathrm{Sp}^{\mathrm{ch}}}
\providecommand{\Coh}{\mathrm{Coh}}
\providecommand{\CoHA}{\mathrm{CoHA}}
\providecommand{\Meff}{\mathcal{M}^{+}_{\mathrm{eff}}}
\providecommand{\Heq}{H^{\bullet}_{\mathrm{eq}}}
\providecommand{\PhiFA}{\Phi^{\mathrm{FA}}}
\providecommand{\C}{\mathbb{C}}
\providecommand{\ZZ}{\mathbb{Z}}
\providecommand{\bP}{\mathbb{P}}
\providecommand{\cO}{\mathcal{O}}
\providecommand{\cF}{\mathcal{F}}
\providecommand{\cZ}{\mathcal{Z}}
\providecommand{\Rep}{\mathrm{Rep}}
\providecommand{\frakg}{\mathfrak{g}}
\providecommand{\fgl}{\mathfrak{gl}}
\providecommand{\Eone}{E_1}
\providecommand{\Etwo}{E_2}

\section*{Positive Geometry, the $D_5$ Character, and $G(X)$}
\label{note:positive-geometry-d5-character}

\subsection*{The Positive Half}
\label{note:positive-half}

Let \(X\) be a CY category or a smooth CY variety in a class where the
effective moduli problem is defined. The positive cone in numerical
\(K\)-theory selects an open and closed union
\[
 \Meff(X)\subset \mathrm{Perf}(X)
\]
of effective objects. On a critical CY$_3$ presentation the
vanishing-cycle sheaf is denoted by \(\phi_W\); in motivic Hall theory
the same symbol records the critical orientation datum. The positive
Hall half is
\[
 Y^{+}(X)
 :=
 \Heq\!\bigl(\Meff(X),\phi_W\bigr),
\]
with convolution along the correspondence of short exact sequences or,
in the derived setting, along the stack of extensions. This is the
Kontsevich and Soibelman cohomological Hall construction in positive
cone form. A quiver with potential \((Q_X,W_X)\) gives a local chart
on \(\Meff(X)\) and localizes the positive geometry.

The layer structure is part of the definition:
\[
Y^{+}(X)
\quad\leadsto\quad
D_{\hbar}^{\wedge}\bigl(Y^{+}(X)\bigr)
\quad\leadsto\quad
\cZ\!\left(\Rep^{\Eone}\bigl(Y^{+}(X)\bigr)\right).
\]
The first term is an associative Hall algebra, the second term is the
completed Drinfeld double after the negative half, Cartan part,
continuous coproduct, Hopf pairing, centre data, and completion are
constructed, and the third term is the derived chiral centre carrying
the half-braiding and \(R\)-matrix. The notation \(G(X)\) is reserved
for the completed double plus centre package when these structures have
been supplied:
\[
 G(X) := D_{\hbar}^{\wedge}\bigl(Y^{+}(X)\bigr)
 \quad\text{together with its centre and completion data.}
\]
For \(X=\C^3\), Schiffmann and Vasserot identify
\[
 \CoHA(\C^3)\cong Y^{+}(\widehat{\fgl}_1).
\]
The full affine Yangian and the \(\mathcal W_{1+\infty}\) realization
enter after Drinfeld doubling and evaluation. The negative modes,
centre, and completed quantum group belong to those later layers.

\subsection*{Equivariance Strata}
\label{note:equivariance-strata}

The symbol \(\Heq\) depends on the geometric stratum:
\[
\begin{array}{lll}
\text{toric CY}_d
& H^\bullet_{(\C^*)^d}
& \text{torus weights on the local critical chart},\\[0.25em]
K3\times E\text{ and reduced compact CY}_3\text{ geometries}
& H^\bullet_{\mathrm{red}}
& \text{translation-reduced obstruction theory and automorphisms},\\[0.25em]
X/G\text{ orbifold}
& H^\bullet(I(X/G))
& \text{inertia-sector grading},\\[0.25em]
\text{smooth projective CY without torus or orbifold chart}
& H^\bullet_{\mathrm{MHS/mot}}
& \text{mixed-Hodge or motivic critical cohomology}.
\end{array}
\]
The toric row recovers the usual quiver and shuffle presentation. The
compact row is the \(K3\times E\) row: the translation direction on the
elliptic factor creates the reduced obstruction theory used in the
Oberdieck and Pandharipande reduced DT scalar. The orbifold row supplies
the CHL inertia sectors. The last row names the required replacement
for torus weights in the absence of a localization presentation; a
complete Hall algebra in that row requires a Joyce and Song or Davison
and Meinhardt motivic critical construction.

\subsection*{The CY to Chiral Boundary}
\label{note:cy-to-chiral-boundary}

For \(d=3\) the CY to chiral construction is the tagged object
\[
 \Phi_3^{(\Sigma_2,C)}(\mathcal C)
 =
 \SpCh_{\Sigma_2,C}\!\left(\PhiFA_3(\mathcal C)\right).
\]
Stage \(1\) produces the \(E_3\) holomorphic factorization algebra.
Stage \(2\) integrates over the complex surface \(\Sigma_2\) and
restricts to the reference curve \(C\), so the output on \(C\) is
\(\Eone\)-chiral. The \(\Etwo\) structure attached to a \(d\geq 3\)
output belongs to
\[
 \cZ\!\left(\Rep^{\Eone}
 \bigl(\Phi_3^{(\Sigma_2,C)}(\mathcal C)\bigr)\right).
\]
For \(Y=K3\times E\) the canonical tag is
\[
 A_Y^{(K3,E)}
 =
 \SpCh_{K3,E}\!\left(\PhiFA_3(D^b\Coh(Y))\right)
 \in \Eone\text{-}\ChirAlg(E).
\]
This chiral object and the Hall positive half \(Y^+(Y)\) are compared
through the Hall to chiral maps, the Drinfeld double, and the centre
data.

\subsection*{The $K3\times E$ Invariants}
\label{note:k3xe-invariants}

Let \(Y=K3\times E\). Four invariants occur on four different
constructions:
\[
\begin{array}{lll}
\kcat(Y) & =0
& \chi(\cO_Y)=\chi(\cO_{K3})\chi(\cO_E)=2\cdot 0,\\[0.25em]
\kch^{\mathrm{Heis}}(Y) & =3
& \text{Heisenberg and Mukai Stage \(2\) specialization on }E,\\[0.25em]
\kBKM(\Delta_5) & =5
& c_1(0)/2=10/2\text{ for the primitive BKM denominator},\\[0.25em]
\kfib(Y) & =24
& \mathrm{rk}\,\widetilde H(K3,\ZZ)=1+22+1.
\end{array}
\]
The compact Hodge supertrace
\[
 \sum_q(-1)^q h^{0,q}(Y)=1-1+1-1=0
\]
agrees with \(\kcat(Y)\) on the total space. The K3 fibre value
\(\chi(\cO_{K3})=2\) is a fibre witness. The total-space invariant
list is \(\{0,3,5,24\}\).

\subsection*{The $D_5$ Character and the BKM Denominator}
\label{note:d5-character}

The primitive BKM denominator and the Oberdieck and Pandharipande
reduced DT scalar use different normalizations. The primitive K3
Borcherds product is \(\Delta_5\), of weight \(5\), and
\[
 \kBKM(\Delta_5)=c_1(0)/2=5.
\]
The unnormalised Igusa square is
\[
 \Phi_{10}^{\mathrm{un}}=\Delta_5^2.
\]
The OP scalar uses the monic Borcherds product
\[
 D_5=64^{-1}\Delta_5,
 \qquad
 \Phi_{10}^{\mathrm{OP}}=D_5^2=4096^{-1}\Delta_5^2.
\]
Hence the reduced DT scalar is
\[
 Z_{\mathrm{DT}}^{\mathrm{red}}(K3\times E)
 =
 -\bigl(\Phi_{10}^{\mathrm{OP}}\bigr)^{-1}
 =
 -D_5^{-2}
 =
 -4096\,\Delta_5^{-2}.
\]
This scalar is a character constraint. The compact Hall source,
primitive Hall bracket, Hopf pairing, centre, and completion are extra
structure. BKM recognition requires the Borcherds denominator identity,
the root multiplicity formula, and finite Hall comparison maps from
\(Y^{+}(K3\times E)\) to the Borcherds positive half.

\subsection*{Six Constructions and One Conditional Target}
\label{note:six-constructions}

The six constructions attached to \(K3\times E\) have distinct inputs:
Schiffmann and Vasserot CoHA, Maulik and Okounkov stable envelopes,
Gritsenko and Nikulin Borcherds lift, Costello and Paquette holographic
boundary algebra, Bridgeland and Smith autoequivalence monodromy, and
Fourier-Jacobi genus escalation from primitive BPS invariants. Their
common target is the completed Hall Drinfeld package
\[
 \mathscr D_{\Delta_5}
 =
 D^\wedge_\hbar\!\bigl(Y^+_{\mathrm{Hall}}(K3\times E)\bigr)
\]
under finite comparison maps preserving the PBW root grading,
Serre kernel, coproduct, nondegenerate pairing, allowed centre, and
pro-Borcherds completion. The common scalar character is
\[
 -\bigl(\Phi_{10}^{\mathrm{OP}}\bigr)^{-1}
 =
 -D_5^{-2}
 =
 -4096\,\Delta_5^{-2}.
\]
Agreement of this scalar with the OP reduced DT series is necessary
for the six-route identification. Construction of the compact critical
CoHA and identification of the completed Drinfeld double are stronger
finite-height obligations.

\subsection*{Proof Obligations}
\label{note:proof-obligations}

\begin{enumerate}
 \item Construct the oriented compact critical CoHA for
 \(K3\times E\), including reduced obstruction theory, vanishing-cycle
 orientation, compact support, and a continuous Hall coproduct.
 \item Prove that the Hall positive half maps to the Borcherds
 positive half with the correct primitive multiplicities
 \(c(4nm-\ell^2)\) and with the OP scalar normalization separated
 from the primitive denominator \(\Delta_5\).
 \item Construct the negative half, Cartan part, nondegenerate Hopf
 pairing, centre data, and pro-Borcherds completion needed for
 \(D^\wedge_\hbar(Y^+(K3\times E))\).
 \item Identify the \(\Eone\)-chiral output
 \(A_Y^{(K3,E)}\) with the Hall package by explicit Hall to chiral
 comparison maps. The \(\Etwo\) structure remains on
 \(\cZ(\Rep^{\Eone}(A_Y^{(K3,E)}))\).
 \item Replace \(H^\bullet_{\mathrm{MHS/mot}}\) by a complete motivic
 critical Hall construction for smooth projective CY$_3$ geometries
 without torus, orbifold, or reduced compact shortcuts.
\end{enumerate}

\subsection*{Primary Sources}
\label{note:primary-sources}

\begin{enumerate}
 \item Kontsevich and Soibelman, 2008: cohomological Hall algebras and
 stability data on effective moduli.
 \item Schiffmann and Vasserot, 2013: the \(\C^3\) positive-half
 identification \(\CoHA(\C^3)\cong Y^{+}(\widehat{\fgl}_1)\).
 \item Oberdieck and Pandharipande, 2016: reduced curve counting on
 \(K3\times E\) and the OP scalar normalization.
 \item Gritsenko and Nikulin, 1995; Borcherds, 1998: the K3
 denominator \(\Delta_5\), Borcherds product, and weight extraction
 \(\kBKM=c_1(0)/2=5\).
 \item Beilinson and Drinfeld, 2004; Maulik and Okounkov, 2012:
 chiral centres, stable-envelope \(R\)-matrices, and the centre
 location of the braided structure.
\end{enumerate}
